%% 第 7 章：复向量空间上的上三角化
%% 本文件由 tools/md2tex.py 自动生成，请勿直接编辑。

\chapter{复向量空间上的上三角化}


\begin{jthm}{上三角化定理}{r:8:1}
若 $V$ 是有限维复向量空间，则任意算子：

\[
T:V\to V
\]

都至少存在一组基，使其矩阵为上三角矩阵。
\end{jthm}

\section{为什么只有复数域下总能上三角化}

证明的起点是代数基本定理。

$T$ 的特征多项式是一个复系数多项式，因此至少有一个复根。也就是说，$T$ 至少有一个特征值与一个一维不变子空间。

之后将该一维不变子空间取商，在商空间上继续寻找特征值。不断重复，便得到一条逐层增长的不变子空间链。这里的证明从略。

\begin{jthm}{舒尔定理}{r:8:2}
在复内积空间下，则总存在至少一组规范正交基使其矩阵为上三角矩阵，这称为舒尔定理
\end{jthm}

\section{上三角矩阵与不变子空间链}

若存在：

\[
\{0\}\subset V_1\subset V_2\subset\cdots\subset V_n=V,
\]

其中：

\[
\dim V_j=j,
\]

且每个 $V_j$ 对 $T$ 不变，那么可以选取基：

\[
(v_1,\ldots,v_n),
\]

使：

\[
V_j=\operatorname{span}(v_1,\ldots,v_j).
\]

于是：

\[
T(v_j)\in\operatorname{span}(v_1,\ldots,v_j).
\]

因此，$T(v_j)$ 的坐标在第 $j$ 行以下都为零。由于矩阵第 $j$ 列记录 $T(v_j)$，矩阵必为上三角形式。

\begin{jinsight}
上三角矩阵的本质，是一条逐层嵌套的不变子空间链。
\end{jinsight}

\begin{jprop}{上三角矩阵的可逆性}{r:8:3}
对于一个算子 $T$ 在某组基下的上三角矩阵，它是可逆的当且仅当全部对角线元素非零

证明思路：若存在对角线 $0$ 元素，则意味着，存在多个向量能够被映射为同一个结果，这意味着映射并非单射，因此零空间非零。
\end{jprop}

\begin{jprop}{上三角矩阵的特征值}{r:8:4}
若 $T$ 的上三角矩阵对角元为：

\[
a_1,\ldots,a_n,
\]

则这些对角元就是 $T$ 的全部特征值，计入代数重数。
\end{jprop}

\begin{jproof}{证明思路}
\[
T-\lambda I
\]

仍然是上三角矩阵，其对角元是：

\[
a_1-\lambda,\ldots,a_n-\lambda.
\]

上三角矩阵可逆，当且仅当所有对角元都非零。

因此，$T-\lambda I$ 不可逆，当且仅当：

\[
\lambda=a_j
\]

对某个 $j$ 成立。
\end{jproof}
